% ============================================================
%  Harald Høffding, PHILOSOPHY AND THEOLOGY. A Historical-Critical
%  Treatise [Philosophie og Theologie. En historisk-kritisk Afhandling].
%  Copenhagen: C. W. Stinck, 1866.  46 pp.
%  TRANSLATION (English)
%
%  Translated from transcription.tex (Danish) in this directory, whose
%  header records how its text was established from the KB scan (every
%  page transcribed at the image and read a second time).  Method:
%  ../../../TRANSLATION-PLAYBOOK.md.  No published translation was
%  consulted (none is known).
%
%  ---- REGISTER ----
%  Scholarly, moderately literal, readable English; American spelling.
%  The treatise is argumentative and quotes heavily (Nielsen, Martensen,
%  Kierkegaard's pseudonyms, Strauß, Stilling, Brandes); quotations are
%  kept as quotations with their page references, and Høffding's long
%  periods are broken only where English will not carry them, never
%  across a page marker or a footnote anchor.
%
%  ---- TERMINOLOGY ----
%  Aligned with the Nielsen, Brøchner and Brandes translations in this
%  repo:  Tro / Viden = faith / knowledge;  Erkjendelse = cognition;
%  Videnskab = science (Wissenschaft), videnskabelig = scientific;
%  ueensartet / eensartet = heterogeneous / homogeneous;  Aand = spirit;
%  Existents = existence;  Inderlighed = inwardness;  Subjectivitet =
%  subjectivity;  Alvor = earnestness;  Aabenbaring = revelation;
%  Forstand = understanding;  Fornuft = reason;  Troesvished = certainty
%  of faith;  indirekte Meddelelse = indirect communication;  Dogmatik =
%  dogmatics (the discipline), Dogmatiken = the Dogmatics (Martensen's
%  book);  Christenheden = Christendom, Christendom = Christianity;
%  Kritikeren = the critic (Rasmus Nielsen).
%  Kierkegaard's works in their standard English titles inside the
%  author's quotation marks; „Uvidenskabelig Efterskrift“, Høffding's
%  short form, = ``Unscientific Postscript'' (abbreviated as he
%  abbreviates it).  Martensen's „Den christelige Dogmatik“ = ``Christian
%  Dogmatics''.  All other titles stay in their original language as
%  printed (Samvittighed og Videnskab, Dogmatiske Oplysninger, undersøgende
%  Anmeldelse, Evangelietroen og Theologien, Grundideernes Logik ...).
%
%  ---- CONVENTIONS ----
%  - Page markers \opage{N} and "% --- p. N ---" at the same joints as in
%    transcription.tex, each with \setcounter{footnote}{0} (notes are
%    numbered 1), 2) ... afresh on every page, as printed); a word the
%    Danish divides across a page turn is kept whole in English.
%  - The same paragraphs, \emph spans (letterspacing in the original),
%    \textit spans (antiqua: Latin and foreign words), footnotes at the
%    same anchors, rules and centre blocks.
%  - Quotation marks „...“ -> ``...''.  Latin tags that the print sets in
%    guillemets stay in guillemets (»credo ut intelligam«); Latin, German
%    and French stay in their language, Greek verbatim.  Bare "p." before
%    page numbers as printed.
%  - Translator's decisions are recorded in "% TR:" comments at the site.
%
%  ---- DEFECTS: DANISH AS PRINTED, ENGLISH CORRECT ----
%  The transcription carries every printer's error.  The translation
%  renders the evident SENSE and logs the divergence in a % TR: comment:
%  p. 9 "videnstabeligt", "Parodox"; p. 16 an unclosed parenthesis;
%  p. 19 „indrekt Meddelelse“; p. 21n "væsenlig"; p. 22 "Objectiviten";
%  p. 24 „den objektive Vished“ in a quotation from the Postscript, which
%  reads Uvished, rendered "uncertainty" (the print, confirmed at the
%  image, has Vished); p. 28 "eu"; p. 29 »contra“; p. 30 "euhver";
%  p. 31 "Pseydonymerne"; p. 41 »Suum cuiqve!»; unbalanced quotation
%  marks supplied in English where the print omits them (pp. 28n, 32).
%  German quoted as printed ("kan", "inthalten", p. 36n).
%
%  ---- HOW THE TRANSLATION WAS CHECKED (2026-09-23) ----
%  (1) Five translators, one per batch (title page and pp. 3--11, 12--20,
%  21--29, 30--38, 39--46), each re-reading its English against the Danish
%  sentence by sentence before returning.  (2) Mechanical audit
%  (tr_audit.py): per printed page the same \opage joints, \emph/\textit
%  spans, footnotes, rules, centre blocks and paragraph breaks as
%  transcription.tex, and an English/Danish word ratio within 0.95--1.60:
%  44 pages and the title page, 0 flagged.  (3) Independent review: three
%  readers who had not translated the pages read every sentence of the
%  Danish against the English.  No omission of substance found; 21
%  corrections applied, among them p. 6 "først" = only, not first; p. 24
%  Vished/Uvished (above) and a dropped "dog"; p. 25 a past
%  counterfactual; p. 29 "vel maa" = may; p. 35 "i endelig Forstand" = in
%  a finite sense, and "Naar" = since; p. 39 "kun komisk" = merely
%  comical; p. 41 "Videnskabens Objective" = what is objective in
%  science; p. 45 "først" = only.  Seams between batches (pp. 20/21,
%  29/30, 38/39) checked by the reviewers.
%
% ============================================================
\documentclass[12pt,a4paper]{book}
\usepackage[utf8]{inputenc}
\usepackage[T1]{fontenc}
\usepackage{libertinus}
\usepackage{libertinust1math}
\usepackage{textalpha}   % Greek (pp. 10, 23)
\usepackage[english]{babel}
\usepackage[protrusion=true,expansion=true]{microtype}
\usepackage{setspace}
\usepackage[margin=1.2in]{geometry}
\usepackage{enumitem}
\usepackage{fancyhdr}
\setlength{\headheight}{28pt}
\addtolength{\topmargin}{-13.5pt}
\usepackage{hyperref}

\hypersetup{
  pdftitle={Philosophy and Theology},
  pdfauthor={Harald Høffding},
  hidelinks
}

\renewcommand{\thefootnote}{\arabic{footnote})}

\pagestyle{fancy}
\fancyhf{}
\fancyhead[LE,RO]{\thepage}
\fancyhead[RE]{\textit{Philosophy and Theology}}
\fancyhead[LO]{\textit{Harald Høffding}}

\setstretch{1.3}

\usepackage{marginnote}
\newcommand{\opage}[1]{\marginnote{\footnotesize\textit{[#1]}}}

\begin{document}
\thispagestyle{empty}

% TR: Danish title page: "Philosophie og Theologie. En historisk-kritisk
% TR: Afhandling af Harald Høffding, cand. theol. Kjøbenhavn. Bog- og
% TR: Papiirhandler C. W. Stinck's Forlag. Thieles Bogtrykkeri. 1866."
\begin{center}
{\Huge Philosophy and Theology.}\par
\rule{4em}{0.4pt}\par\bigskip
{\large A Historical-Critical Treatise}\par\medskip
by\par\medskip
{\large Harald Høffding,}\par
\textit{cand. theol.}\par\bigskip
\rule{6em}{0.8pt}\par\bigskip
Copenhagen.\par
% TR: publisher's and printer's imprint kept as printed (publisher: bookseller
% TR: and stationer C. W. Stinck; printer: Thiele); the second initial "W." is
% TR: the transcription's reading (alternatively V.).
Bog- og Papiirhandler C. W. Stinck's Forlag.\par
{\small Thieles Bogtrykkeri.}\par
1866.
\end{center}
\clearpage

% --- p. 3 ---
\opage{3}\setcounter{footnote}{0}%
When a debate concerning an essential point of contention
returns again and again, when the discussion is taken up
anew again and again, then it may seem trivial, to an outward
view, always to be quarreling in this way about the same thing. But
it will also be so only for the outward, superficial view.
In reality the discussion must always return to the same
field of battle with richer equipment, furnished with the new
means that the advancing development has brought forth since
the last time. This finds its application when one considers
the different conditions and presuppositions under which the
disputes between philosophy and theology have been
conducted at different times, and particularly when we look to
the peculiar way in which the dispute has developed
in \emph{our} literature.

Whenever, therefore, the matter is stirred up anew, it is
natural to turn one's gaze back to the earlier history of the
matter, in order to see whether the new contribution has really
understood it and made right use of it --- for that is the first
condition for getting any further. But here things may then turn out
so strangely that the dispute is really no longer conducted to settle
the matter itself, but to settle who actually \emph{has}
won in the decisive dispute. For in a dispute of ideas things
do not go as in the battles of the bodily world, where the
decision is something one can take hold of and touch.

% --- p. 4 ---
\opage{4}\setcounter{footnote}{0}%
% TR: the title of Larsen's book, "Conscience and Science", is kept in Danish
% TR: throughout, as the brief directs for titles other than Kierkegaard's and Martensen's.
When I read the new contribution\footnote{``Samvittighed og Videnskab'', by A. C. Larsen, \textit{cand. theol.}} to the dispute mentioned above,
which, from what I have heard, is regarded by many as having
no small significance, particularly as an attack on the
view of the relation between theology and philosophy, and
between faith and knowledge, maintained by Prof. R. Nielsen,
I got an example of this. Although the author seems to regret that ``since the
dispute that was conducted in its time over Martensen's Dogmatics, all
polemic about the relation of faith and knowledge to each other has
fallen completely silent'', he has nevertheless not done what was in his
power to carry the matter further in the right way: namely,
to give himself a clear account of how matters stand. This
has harmed his work all the more, as the historical
course of development points fairly clearly toward the view that
must be the result of the long, entangled dispute. A historical
presentation of this course of development will be able at once
to show us which the victorious view must be, and
to furnish a critique of the work named.

\begin{center}\rule{4em}{0.4pt}\end{center}

Every student, and particularly one who has sat
on the pupils' bench in the philosophical as well as in the theological
lecture hall, must feel called upon to make clear to himself
what it really is that is disputed here, and how he himself
must take his stand in this matter. This is by no means easy for the
young man, who must essentially take the part of a learner; for here
authority stands against authority, and he who now perhaps sees shaken
what he took to stand fast as a rock easily becomes
disheartened and wavering. But precisely by this he is then
prevented from one thing: what the Latins call \textit{jurare in verba magistri} ---
he cannot, after all, swear by two authorities at once! So
he then sees the necessity of taking his side according to his own
best conviction.

% --- p. 5 ---
\opage{5}\setcounter{footnote}{0}%
What now leads me to come forward publicly with my
view of the matter, instead of keeping it to myself,
is that it seems fitting to me that the dispute should also be fought out
among the younger men, after the main battle was fought in its time by the
masters themselves.

In this respect I have, to be sure, a predecessor in the author
of the work cited above, which has given me occasion for
this little treatise.

\begin{center}\rule{4em}{0.4pt}\end{center}

% --- p. 6 ---
\opage{6}\setcounter{footnote}{0}%
Before I proceed to my real task: to show
how the view that in my conviction is the
only tenable one emerges from the whole course of the dispute, especially
by reason of the turn this has taken in our literature through
S. Kierkegaard's appearance, --- I shall briefly indicate why
I must regard Mr. Larsen's contribution as misconceived.

The author of ``Samvittighed og Videnskab'' naturally acknowledges
a difference between faith and knowledge; but he will not
concede that they are absolutely heterogeneous. He regards it as
impossible to hold fast the conviction of the reality of modern knowledge
and the assurance of the reality of Christianity, except
by the help of scientific theology, where faith and knowledge
can negotiate with each other (p. 60).

But in this way he only reconciles the two opposites
by confounding them. He makes science
religious in order to be able to make religion scientific.

When I say that the author makes science religious, I
mean by this that he refers problems that belong to
philosophy to religion. --- Whence do we know that external
sense perception has reality? This question, in the author's
opinion, cannot be decided ``by way of \emph{knowledge}''. The external world
comes to be real for me only through \emph{faith} (p. 13--14), and
this faith is in principle one with that in which God exists
for me (p. 14). This is an equivocation. Philosophy
is really able ``by way of knowledge'' to solve the problem of objectification\footnote{Reference may be made here to ``Grundideernes Logik'' I, p. 172--177.}.
And insofar as one can speak of a philosophical
% --- p. 7 ---
\opage{7}\setcounter{footnote}{0}%
\textit{sapere aude}, a faith in the reality of thought, then
this faith is in principle different from the religious one, even
if it can be supported by it. This can be seen from the relation
between dogmatic and critical philosophy: ``die kritische
Philosophie begreift, was die dogmatische geglaubt hatte''\footnote{Kuno Fischer, Geschichte der neuern Philosophie I, p. 241.},
and what dogmatic philosophy had believed in was precisely
the reality of cognition. --- What is the relation between spirit
and nature? About this, in the author's opinion, one can learn nothing
in science: ``for natural science as such has
no concept of nature; as such it knows nothing of
nature's relation to the Absolute. The powers between which
the struggle is waged, without reconciliation ever being possible, are
pantheism and Christianity'' (p. 38). This is simply
incorrect. For it is overlooked that there is something called
philosophy of nature: here one has a concept of nature, even if physics
otherwise has none; and here the relation between idea and nature is shown.
Here one can learn how the highest cause works in and
through all the finite causes. The philosophy of nature will neither
believe nor comprehend the dogma of creation; but it will comprehend the relation
between spirit and nature as a fundamental relation. If
one wants to call this pantheism, that is an equivocal expression;
for the philosophy of nature is not religion but a scientific
grounding of the phenomena of nature.

If the author thus\footnote{One cannot help thinking here of Hegel's words: ``Sie ignoriren die Philosophie, aber nur um in ihren willkürlichen Raisonnements nicht genirt zu werden''. Geschichte der Philosophie I, p. 80.} overlooks philosophy, he is
all the better friends with physics. What an \textit{entente cordiale}
there is for him between ``the real laws of nature'' on the
one side and the dogma of creation and the miracles on the other
side! For the dogma of creation implies that nature is a real
``something'': so there are also ``real laws'' and a ``real
natural science'' (p. 33)! But what, I wonder, will natural science
% --- p. 8 ---
\opage{8}\setcounter{footnote}{0}%
say when one lets the dogma speak its piece: ``it is something ---
and yet nothing for the Almighty!'' I do not believe it helps
to explain to the natural scientist that now we pass over into another
science, where one has this ``other way of seeing
things''. --- And now the miracle! Yes, says the author, it ``precisely
confirms the law of nature; if one posits the exception, one must also
posit the rule'' (p. 37). A splendid piece of reasoning --- especially
for thieves and robbers, who can thus earn themselves merit
by enforcing the validity of the law, if in a rather
``indirect'' way!

The author is right that it is not Christianity and natural science
``that one must choose between''; but he has no right
to put ``theology'' in place of ``Christianity'' (which
he does p. 37--38); for this changes the whole matter.
For the dispute between theology and natural science is a
dispute between two sciences.

But it is not so strange after all that the author --- consciously
or unconsciously --- identifies Christianity and theology. For
he has only made science religious in order to make religion
scientific.

When the author wants to make religion an object of science,
it is by no means his opinion that one can ``just like that,
without further ado'' comprehend the religious truths. But he
thinks that ``it is something quite different when I, after having
gone out beyond myself in faith and after having posited God
as cause, then afterwards infer back from my God-consciousness, as effect,
to the constitution of God's being'' (p. 20).
By all this going out and in one then attains ``a
science that essentially wants to be nothing other than the religious
life's own clear and transparent consciousness of
itself'' (p. 54).

But, one must ask, does proving and grounding then go
more easily after one has become a believer than before?
Let it be well noted that the question is not: whether the
% --- p. 9 ---
\opage{9}\setcounter{footnote}{0}%
religious consciousness is convinced of the truth of that
in which it believes, but: whether this conviction can become
a scientific insight and comprehension?

Now Johannes Climacus says that the Christian truth
is the paradox! This hindrance to ``the religious life's
clear and transparent consciousness of itself'' must then first
be cleared out of the way. And that goes easily enough: ``It is one thing to express oneself with
lyrical-passionate energy in a polemical situation;
% TR: Danish prints "videnstabeligt" (misprint for "videnskabeligt"); rendered by its sense.
it is something quite different to speak positively-scientifically'' (p. 49).

No one, surely, disagrees with the author in this; but when he
concludes from it that Joh. Climacus's proposition can hold only
when it is the humorist who has it in his mouth ``in a polemical
situation'', then this is\footnote{How, I wonder, has the author understood the essence and significance of ``indirect communication'', when he can conclude in this way!} a strange mixture of thoughtlessness
and sophistry. The believer's relation to the object of faith
in faith and hope is a direct relation, in which there
% TR: Danish prints "Parodox" (misprint for "Paradox"); rendered by its sense.
is no paradox; but as soon as the believer wants to have objective
(i.e.\ scientific)\footnote{What right has the author to define objective certainty as ``pure logical certainty of state?''} certainty, then the object of faith shows itself
--- according to Joh. Climacus --- as the absurd. But, says
the author, Joh. Climacus himself says, after all, that ``objective indifference''
gets to know nothing at all about Christianity; how
then can one get to know objectively that it is the absurd?
The answer is not so difficult after all: for getting to
know that it is the absurd, and getting to \emph{know}\footnote{Scientifically, of course. --- Examples in Joh. Climacus p. 160 f. 166 f.} nothing at all, are one and the same. The paradox must, before
the scientific tribunal, be declared a \textit{non-ens}. --- Nonetheless
the author regards Joh. Climacus as happily and well
placed in ``a place of his own'' outside the ``positive-scientific''.
It is only, as it were, for good measure that the poor
% --- p. 10 ---
\opage{10}\setcounter{footnote}{0}%
humorist must hear that he does not have ``positive-scientific''
seriousness enough to distinguish between θεός and ὁ θεός (p. 53); and
then at the end he receives the lesson that one must after all
have the objective along too. And from the fact that faith cannot
do without ``the objective'' it is then concluded\footnote{In the same cunning way the author earlier put the word ``theology'' in place of the word ``Christianity''.} quite directly: that neither
can it do without ``scientificity'' (p. 54)!

But enough: the author has finished with Johannes
Climacus; he now establishes dogma as a fact
for the religious consciousness, and engages the understanding (by
the author's own declaration, p. 52, the same\footnote{This, incidentally, does not matter, since by ``understanding'' the author understands something altogether formal (p. 50--51). An understanding so ``reduced'' can be made into anything whatever. The author seems not to know that there is something called philosophical thinking, concept, cognition of ideas. He really makes the task far too easy for himself!} understanding that
the believer had before he became a believer) to set to work
(p. 50--52; 71--74). But here one must regret that the author
has overlooked precisely what matters. He simply sets
religion alongside nature, without thinking
of the essential difference between the two, regarded as objects
of science\footnote{This is all the less excusable, since this parallel has been contested in a separate treatise: ``Er Religionsforskerens Forhold til den saakaldte Aabenbaring det samme som Naturforskerens til Naturen?'' By Mag. Stilling. 1853.}. Of the religious fact
holds what Feuerbach says: ``Thatsache ist kein Grund''.
In the facts of nature, on the other hand, reason can find its
own essence and so ground them. ``Die Vernunft ist
die Hebamme der Natur''. And as it goes with the category
``fact'', so it goes with the related category
``presupposition''. A scientific presupposition must be able wholly and
fully to pass through the fire-test of doubt; this holds
% --- p. 11 ---
also for what is the object of ``the immediate sense \opage{11}\setcounter{footnote}{0}impressions'':
it is not recognized scientifically unless it can be grounded.
But the author cannot go along with that: ``however justified
one may be in demanding doubt, one thing one must nevertheless
not demand, that the religious consciousness should doubt the reality
of its immediate being'' (p. 58.). Yes indeed, that one may and
must demand, if science is to come of it,
and in this respect theology diverges very greatly in its
``demands'' from natural science.

We get a practical test of the author's science in his
3rd treatise, where he wants to show us how God stands
``in real reciprocal relation'' with time (p. 46). Can he now
really develop this reciprocal relation between the eternal, personal,
omniscient God and temporal, successive development, so
that it becomes a scientific concept? He seems, however, only to
show that it ``must be thought'' --- but he does not think it. ---

When the author has thus given neither science nor
faith what is due to each of them, the reconciliation cannot
be lasting either. And that this is so we know not
merely \textit{a priori}, but --- as what follows shows --- also
\textit{a posteriori}.

\begin{center}\rule{4em}{0.4pt}\end{center}

% ==== batch 2: printed pp. 12--20 (PDF 16--24) ====
% p. 12 opens a new section (after the short rule closing p. 11); no heading
% is printed.
% --- p. 12 ---
\opage{12}\setcounter{footnote}{0}%
In the Middle Ages theology, as scholasticism, was the dominant, indeed the only
science. ``Dogma is true (for the Church teaches it!); therefore it also
accords with reason; therefore it can be proved.'' So one reasoned, and then
summoned all the powers of the understanding to build up the great scholastic
systems, which have aptly been compared with the Gothic cathedrals. And it
was a pious work; just as the pious master builders on the Rhine raised spires
pointing toward heaven, where often, high up, where no human eye ever reached,
ornaments were placed to the glory of God, so too the scholastic systems, with
their whole apparatus of syllogisms and distinctions, were erected
\textit{in gloriam Dei}. The scholastics fought with refractory thoughts in the
same way as the crusaders against the Saracens, and it went ill with both if
they did not bow before faith and the Church.

% TR: "Thomas Aqvinas", "Vilhelm Durandus", "Vilhelm Occam" (pp. 12--13) are
% the 1866 Danish forms of the names; given in their usual English forms.
With the medieval Church scholasticism also fell. Its fall was brought about
both by inner self-dissolution and by attack from without. The presupposition
of scholasticism was the unity of faith and knowledge: but now Duns Scotus, by
making the will the highest principle in opposition to Thomas Aquinas, had
already distinguished between the theoretical and the practical, and William
Durandus (in the first half of the 14th century) had recognized that there
were dogmas which could not and should not be proved, but believed as revealed
truths. The presupposition of scholasticism was, further, the unity of
thinking and being:
% --- p. 13 ---
\opage{13}\setcounter{footnote}{0}%
but now, through William of Occam, nominalism (\textit{universalia sunt nomina,
flatus vocis}) had become dominant in a sharpened and renewed form. The
attacks from without came partly from the side of thought, now set free, which
wanted to tear down the prison in which it had so long been tortured, and
partly from the existential, religious standpoint, when Luther broke with the
traditional dogma and all the »\textit{traditiones humanæ}« in which it had
cocooned itself.

But the new philosophical inquiry did not yet dare to attack the doctrine of
the Church itself. In order, then, to be able to maintain the results of their
own cognition unchallenged by the authority of the Church, the philosophers set
up the separation between \textit{duæ veritates, philosophica et theologica},
by the proposition that something could be true in theology, from the
standpoint of dogma, and yet false, or at least unprovable, from the
standpoint of philosophy, of reason. As men who asserted this separation may be
named Pietro Pomponazzo (at the end of the 15th and the beginning of the 16th
century), who concerned himself especially with the doctrine of immortality,
and Vanini (at the end of the 16th and the beginning of the 17th century), whom
the distinction in question did not save from the stake.

That this distinction was sophistical and untenable is easy to see. For dogma,
as the Church laid it before its faithful for obedient acceptance, was not the
evangelical truth with its practical-religious character, but was a
theologumenon with ecclesiastical sanction. ``Scientific and religious
cognition had at that time grown together so ecclesiastically that they could
not be separated. When the philosophers, then, in order to preserve the
scientific character of their work, declared that in all matters of faith they
were willing to submit to the authority of the Church, the Church was
certainly justified in calling the sincerity of such a concession into doubt:
how could anyone hold fast to a truth of reason and yet acknowledge a
theological truth, when the two stand in
% --- p. 14 ---
\opage{14}\setcounter{footnote}{0}%
decisive conflict with each other, and yet both were to count as
scientific?''\footnote{R. Nielsen: Grundideernes Logik p. XXII.}

That on these terms the conflict could not lead to reconciliation is then
evident. It was continued, accordingly, in the Protestant as well as in the
Catholic Church. Beside philosophy there now stood natural science, which has
a continuous history only from the 16th century on.

If at the end of the last century the question had been raised who had won in
the conflict, theology or philosophy, the answer could hardly have been in
doubt. For theology, under the influence of the spirit of the age, had let one
miracle fall after another, in the world of spirit as in that of nature. The
theology of the philosophical, naturalistic 18th century was no longer any
theology at all.

But the general spiritual upsurge at the beginning of this century showed
itself in the field of theology as well. It was above all philosophy that now
became of importance. Rationalism had been rejected and the principle of faith
once more embraced; but in the newer philosophy one saw a possibility of
reconciling this principle of faith with the principle of knowledge. Now, then,
there came a time of flowering for theology, which, for those who
cultivated it, was full of enthusiasm. ``Never had so many from the so-called `better classes' been
seen to choose this study as in that\footnote{The reference is to the years
1824--33 in Germany.} period.... Just as some years earlier the flower of
German youth had placed itself under the banners of war to save the fatherland
from political servitude, so in the years of peace that followed those most
richly gifted in spiritual respects regarded it as the highest honor and the
highest gain to be allowed to serve the Church and to offer it their best
powers.... One knew that the struggle was for a noble goal; one believed in the
possibility of a final victory, in the possibility of a reconciliation of faith
and
% --- p. 15 ---
\opage{15}\setcounter{footnote}{0}%
knowledge, even if one did not see it fully accomplished;... one did not yet
regard it as bearing a foreign yoke when theology went hand in hand with
philosophy\footnote{Hagenbach: ``Om Aftagelsen af det theol. Studium'', in
Clausens Tidsskrift for udenl. theol. Litt. 1858, p. 13, 21.}.''

It was above all the Hegelian philosophy that was rich in promises of a
reconciliation of faith and knowledge. Its content was the Absolute --- but
the content of Christianity, too, is surely the absolute truth\footnote{It is
said of the religious spirit that its content is ``als religiös wesentlich
spekulativ'', and of the content of philosophy that it is ``spekulativ und
damit religiös''. Hegel, Encyclopädie der philos. Wissenschaften. Berlin 1845.
p. 513.}. It possessed besides, in its dialectical method, the capacity to
carry through the mediation between the contending principles. How rich, then,
were the presuppositions here for a speculative theology!

% TR: "Forestillingen" (Hegel's Vorstellung, quoted just before) = "representation".
But did Hegel, then, make no distinction at all between faith and knowledge?
Yes, he did, but an ambiguous one, which was therefore afterwards interpreted
\textit{in utramque partem}. ``Worauf es ganz allein ankommt, ist der
Unterschied der Formen des spekulativen Denkens von den Formen der Vorstellung
und des reflectirenden Verstandes''\footnote{Op. cit. p. 512.}. This was now
understood in two ways. Some said: ``Representation is a lower, consequently an
inadequate, and so a relatively untrue form, which must therefore be dissolved
by thought.'' But others held that one could well keep the Hegelian
determination of faith as an immediate knowledge, ``only that we do not then
understand by the immediate the lower, the imperfect, which is to be abolished
by philosophy as the perfect, but the original, the primitive knowledge, which
is the foundation of speculation\footnote{Martensen, Christian Dogmatics. 1849.
p. 13.}.''

This latter view then became the basis of a new tendency in theological
dogmatics, which also found representatives at our university. It was above
all the present Bishop
% --- p. 16 ---
\opage{16}\setcounter{footnote}{0}%
\textit{Dr.} Martensen who, by his\footnote{\textit{De Autonomia
conscientiæ sui humanæ in theologiam dogmaticam nostri temporis introducta.
Hauniæ 1837.}} dissertation for the licentiate degree, paved the way in this
country for this ``new era.'' He would by no means grant philosophy primacy over
theology, but would instead assert a principle of its own for theological
speculation; while philosophical thinking is autonomous, theological thinking
is theonomous (\textit{dogma de deo, auctore conscientiæ sui humanæ,
verum esse theoriæ cognitionis fundamentum, qvo efficitur nexus necessarius
inter scientiam theologicam et auctoritatem dogmatis, vinculum indissolubile
scientiæ et conscientiæ}\footnote{Op. cit. p. 135.}).
% TR: the print never closes the parenthesis opened before "dogma de deo";
% the English supplies the closing parenthesis after the footnote mark.
The concept of mediation was retained, but it was demanded that ``mediation
should be sought from the standpoint of believing cognition''\footnote{Dogmatiske
Oplysninger, by H. Martensen. 1850. \textit{pag.} 66--69. --- ``The center of
Christianity, the doctrine of the Incarnation, shows precisely that Christian
metaphysics cannot rest in an Either---Or'', as does ``the metaphysics of the
Jewish religion.'' Tidsskrift for Litteratur og Kritik 1839, p. 458.}.

Here, then, a dogmatic principle had been laid down which could supposedly
maintain the scientific character of theology by uniting in itself the
principle of faith with that of knowledge. But what our literature saw, even in
the first years, of productions in this direction was only isolated\footnote{Here
may be named, e.g., ``Mester Ekkardt'', the excellent ``Bidrag til at oplyse
Middelalderens Mystik,'' by \textit{Dr.} Martensen (1840).} attempts, not any
thoroughgoing scientific presentation of the doctrine of faith according to
the stated principle. ``The system was not yet finished'' --- but it was
expected, and all looked forward with great hopes to the further development
of the ``new era.''

Before we continue the historical account further, we must first call to mind
the tendency which --- likewise springing from the Hegelian philosophy --- went
the opposite way from the ``believing'' speculation. It is Strauß and
% --- p. 17 ---
\opage{17}\setcounter{footnote}{0}%
Feuerbach who are the chief men here. If the answer that knowledge gave, when
its testimony was demanded after it had been set face to face with revelation,
was for the former tendency unconditionally affirmative, here it became just
as unconditionally negative. Negative criticism thus sets itself against
positive dogmatics, \textit{contra} against \textit{pro}.

The opposition, then, is here as strong as possible. But should there not
after all be something common to the two contending parties, something on
which, in spite of all disagreement, they are agreed?

What both parties were agreed on was that the conflict was to be fought out on
the ground of science and with the weapons of science. This had been stated
from the very outset. When Strauß's ``Leben Jesu'' had appeared, the Prussian
Ministry of Education asked the famous Neander for a statement on whether it
would be right to issue a ban on this book, which was causing so great a stir.
In this statement it says: ``It is the calling of the theological teachers to
work toward this, that the book, which has issued from a tendency opposed to
the spirit of the Christian Church, may come to contribute to the advancement
of science and thereby to promote the Church's own interest as well. But this
goal can be reached only if this struggle is waged with the weapons of science
alone''\footnote{Luther, on the other hand, on a similar occasion asked the
Saxon princes to insist ``daß alleine mit dem Worte Gottes in diesen Sachen
gehandelt werde, wie den Christen gebühret.''}.

But then does the spectacle from the end of the Middle Ages not seem to repeat
itself? Here are »\textit{duæ veritates}« --- and both, after all, claim to be
scientific.

On closer consideration, however, this proves not to be so. We shall here only
ask provisionally: what is it that Strauß and Feuerbach attack, religion or
theology, --- and what does it mean when they declare that they stand on the
standpoint of science? --- The first question is soon answered:
% --- p. 18 ---
\opage{18}\setcounter{footnote}{0}%
% TR: Strauß's ``Dogmatik'' = Die christliche Glaubenslehre (1840--41); the
% author's short title is kept as printed.
Strauß, besides ``Leben Jesu'', in which he dissolves the evangelical history
into myths, has of course also written a ``Dogmatik'', in which the dogmatic
propositions are dissolved one after another into nothing; and Feuerbach does
distinguish between religion and theology as between the immediate,
involuntary intuition and conscious reflection, but he turns his weapons
against both. What is combated, then, is not merely a truth that claims to be
scientific; it is religious truth, the truth of life itself, that is combated
here. And as regards the second point, negative criticism holds fast to the
fixed and unshakable laws in the world of spirit and of nature, as these
prove themselves to human consciousness through science. What rational science
cannot check and ground is for this standpoint a \textit{non-ens}.

As surely as religion is one thing and theology another, so surely must the
Strauß-Feuerbach criticism be viewed from a double point of view. Over against
theology it becomes a conflict between knowledge and knowledge; over against
religion it becomes a conflict between knowledge and faith.

``The conflict with Strauß is transformed at once into a conflict between
religion and science. It is the critics' merit to have given the matter so
strong a push that precisely thereby it has been driven to this decisive
turning point. The conflict of the faith of the Gospel with the Straußian
criticism is a conflict of competence, a conflict of sovereignty, a conflict
over the supreme power, the highest authority in the world of
spirits''\footnote{``Evangelietroen og Theologien''. Tolv Forelæsninger, by
R. Nielsen. 1850. \textit{pag.} 85.}.

One sees, then, how matters stood. On the one side stood dogmatics with the
declaration that knowledge agreed with faith, on the other side criticism with
the equally categorical declaration that knowledge conflicted with and
consequently dissolved faith. It really seemed to be high time that
% --- p. 19 ---
\opage{19}\setcounter{footnote}{0}%
attention should be directed to the ``conflict of competence'' between faith
and knowledge; for on this --- and in this dogmatics and criticism\footnote{Cf.
Martensen's Dogmatics § 17, Remark.} were equally dogmatic --- one got the
requisite enlightenment from neither of the contending parties.

At this point we now turn to the Kierkegaardian writings. If anywhere, then in
these writings everything is summoned to establish and maintain the concept of
faith in its distinctiveness.

We shall here consider Kierkegaard only from the point of view just stated. His
great importance of course holds for the whole age, and his appearance was
indeed determined by the spiritual condition of the time as a whole. He is a
rare example of how a gift and a genius such as falls to the lot of only few
human beings can be exclusively taken possession of (as it were, taken captive) by
the ethical-religious aim: himself existing in the truth, to be for others a
teacher in existing in the truth. But what is peculiar to him lies in his
method --- his maieutics, to recall Socrates. He saw his contemporaries as
caught in an enormous illusion of the senses, namely this, that ``we are all
Christians''; an illusion of the senses can never be removed directly ---
% TR: the print has "indrekt Meddelelse", a printer's error for "indirekte"
% (a reader has pencilled in the missing "i"); rendered "indirect communication".
hence ``indirect communication''! If the sickness rests on taking the phenomenon
for the essence, then the remedy must consist in setting discord between
essence and phenomenon, so that the deception may cease and the essence be
recognized as essence.

It has often been objected to this indirect communication that it was too
``artificial'', too reflective: ``Christianity, religion in general, is surely
conditioned precisely by the immediate relation to the divine!'' To this
Kierkegaard replies: one is mistaken in thinking that it is the remedy that
introduces reflection; no, reflection --- that is precisely the misfortune ---
is there beforehand, but only as half-reflection,
% --- p. 20 ---
\opage{20}\setcounter{footnote}{0}%
% TR: "Christenheden" = Christendom (the Christian world), distinct from
% "Christendom" = Christianity.
false reflection. ``The situation is Christendom, which is a determination of
reflection. In Christendom --- to become a Christian is either to become what
one is (the inwardness of reflection, or the reflection of inward deepening),
or it is first to be torn out of a delusion, which again is a determination of
reflection''\footnote{``The Point of View for My Work as an Author'',
\textit{pag.} 33.}.

% TR: "Uvidenskabelig Efterskrift" (the author's short form of the title) =
% ``Unscientific Postscript''.
On this, then, rests the explanation of ``the ambiguity or duplexity in the
whole authorship, whether the author is an aesthetic or a religious author''.
In this connection it is now neither the aesthetic nor the religious writings
to which we would chiefly refer, but what Kierkegaard himself has designated
as ``the turning point in the whole authorship'': ``Unscientific Postscript'',
by Johannes Climacus.

But what we must no doubt settle first is whether Johannes Climacus really
means what he says. ``He is a humorist, after all, playing and toying with
everything in the most frivolous way!'' Let us hear, then, how he himself
conceived his task when he resolved to become an author. ``Out of love for
men, out of despair over my embarrassing position of having accomplished
nothing and of being able to make nothing easier than it has been made, out of
genuine interest in those who make everything easy, I then conceived it as my
task: everywhere to create difficulties.'' ``If I now dared proclaim loudly
that I had come into the world for, and been called to, counteracting
speculation, that this was my judging work, while my prophetic work was to
foretell a matchless future, as to which men, because I was loud-voiced and
called, could safely rely on what I said --- then there would no doubt be many
a one who, failing to regard the whole thing as a fantastical reminiscence in
a fool's head, would take it in
% TR: "Alvor" = earnestness (here and on p. 21).
earnest''\footnote{How he then came to ``a closer understanding of his own
notion, that he must try to do something difficult'', must be read in his own
words (p. 174 \textit{ff.}).}. Earnestness is not
% TR: the sentence straddles pp. 20/21: "Alvor er ikke [p. 21] saaledes til at
% tage og føle paa." Rendered here as "Earnestness is not"; p. 21 continues
% "something one can take hold of and touch in that way".
% batch 2 ends here: p. 20 ends on the complete word "ikke"; p. 21 continues the sentence.
% --- p. 21 ---
% Batch 3: printed pp. 21--29. p. 21 begins mid-sentence.
\opage{21}\setcounter{footnote}{0}%
% TR: the sentence straddles pp. 20--21 ("Alvor er ikke | saaledes til at tage og føle paa"); p. 21's part is rendered "something one can take hold of and touch in that way", completing "Earnestness is not" (p. 20).
something one can take hold of and touch in that way. Earnestness is the essential relation to the essential, but not directly commensurable with its expression. On the contrary: earnestness is inwardness, and ``the form of opposition is precisely the dynamometer of inwardness.''

% TR: "Misviisning" (a compass's deviation) rendered "error"; "Retning" = "tendency" throughout.
For Johannes Climacus, the error of speculation lay in the whole tendency of the age: ``that, on account of so much knowledge, people had altogether forgotten what it is to exist and what inwardness signifies.'' Therefore he wanted to ``make difficulties'' --- by laying stress on existence. When the existing human being asks: ``What is truth?'' what must be meant is the essential truth, the truth of life, ``the truth that essentially\footnote{%
% TR: Danish "væsenlig" (misprint, t missing) rendered "essential"; "Erkjenden" (the verbal noun) = "cognition". Kierkegaard's "Uvidenskabelig Efterskrift" is given by the author's own shortened title as "Unscientific Postscript" throughout.
``Only ethical and ethical-religious cognition is therefore essential cognition. But all ethical and all ethical-religious cognition is essentially related to the fact that the one who cognizes exists.'' Unscientific Postscript p. 145.} relates itself to existence.'' Truth in general is the agreement, the unity, of subjectivity and objectivity; it thus contains, even viewed purely abstractly, a doubling. But this doubling returns with intensified emphasis when it is not \textit{in abstracto} but existentially that the question of truth is asked. Here the two-sidedness becomes a disjunction. The two ways of cognition go in opposite directions: the one turns toward the objective and leads --- precisely in order to bring out the objective as objectively as possible --- away from the subjective; the other leads to subjectivity, to the appropriation of the truth in inwardness.

The thesis of the subjective, existing thinker then becomes: subjectivity is the truth. There is a difference of sphere between scientific and existential truth. Here no mediation is possible; for by the subject-object of mediation one comes back to abstraction, whereas the question was precisely ``how an existing subject \textit{in concreto} relates itself to the truth.''

% --- p. 22 ---
\opage{22}\setcounter{footnote}{0}%
``But,'' --- it is said --- ``it is surely highly one-sided to say that subjectivity is the truth; subjectivity must, after all, have taken up the objective into itself in order to be the truth.'' This objection rests partly on an ``imperfect and undialectical distinction'' between subjective and objective; partly it has not rightly understood the thesis stated. Let us hear Johannes Climacus. ``While objective thinking is indifferent to the thinking subject and his existence, the subjective thinker, as existing, is essentially interested in his own thinking, is existing in it. Therefore his thinking has another kind of reflection, namely that of inwardness, of possession, by which it belongs to the subject and to no one else . . . The reflection of inwardness is the subjective thinker's double reflection''\footnote{Unscientific Postscript. p. 50---51.}. There is thus not merely a relation of reflection between the objective and thinking, but also between the thought objective and the thinker's existence. But this latter reflection is not logical but practical; it is a \textit{salto mortale}, to recall Jakobi, a ``dialectical deed and act''\footnote{Stilling: ``Om den indbildte Forsoning af Tro og --- Viden''. 1850. p. 30---32.}.

What has been developed so far states the conditions for every existing person's relation to the truth: it is ``the Socratic wisdom, whose immortal merit is precisely to have heeded the essential significance of the fact that the one who cognizes is existing, for which reason he, in his ignorance, was in the truth in the highest sense within paganism''\footnote{Unscientific Postscript. p. 150.}.

If, then, there is to be a standpoint essentially different from the Socratic, this difference must rest on the closer determination of the truth.

The proposition ``Subjectivity is the truth'' contains an opposition to
% TR: Danish "Objectiviten" (misprint for "Objectiviteten") rendered "objectivity".
objectivity. This opposition must then also find its expression in the determination of the truth. While
% --- p. 23 ---
\opage{23}\setcounter{footnote}{0}%
objective thinking seeks certainty in objective grounds, the truth, seen from the subjective standpoint, that of existence, is objective uncertainty. The grounds that can be in question here are merely subjective, existential.

``When subjectivity, inwardness, is the truth, then the truth, objectively determined, is the paradox; and the fact that the truth is objectively the paradox shows precisely that subjectivity is the truth''\footnote{Op. cit. p. 151.}.

The paradox is present where the eternal and unconditioned confronts the temporal and conditioned. For Socrates, too, there was a paradox. ``Yet the eternal essential truth is itself by no means the paradox, but is so by relating itself to an existing person''. For that the infinite relates itself to the finite is, to be sure, a contradiction from the standpoint of existence, one that is overcome only by a \textit{salto mortale}; but from the standpoint of science, of comprehension, this contradiction is mediated\footnote{``Not by a leap but by a transition corresponding to the peculiar character of the functions must comprehending subjectivity in each case move from the finite to the infinite''. Grundideernes Logik p. 230.}.

But if the eternal truth and existing in temporality are now put together in the truth itself, then the truth itself becomes paradoxical --- and then we have the paradox κὰτ᾽ ἐξοχῆν, without mediation being possible.

% TR: the Genesis quotation is translated as Høffding prints it; its "where art thou?" is not part of Gen. 22:1.
Here, then, we stand at revelation. The eternal, personal God speaks to man in time: ``And God tempted Abraham, and said: Abraham, Abraham, where art thou? but Abraham answered: here am I'' (Gen. 22, 1). God himself is born and lives and dies as a human being in time: ``The Word was with God, and the Word was God. And the Word was made flesh, and dwelt among us'' (John 1, 1. 14). Here all historical and speculative knowledge is of no help; the highest that, along the objective way,
% --- p. 24 ---
\opage{24}\setcounter{footnote}{0}%
% TR: the print (checked on the scan) reads "den objektive Vished (paa det sokratiske Standpunkt)"; Kierkegaard's Postscript, which Høffding is quoting, reads "istedetfor den objektive Uvished" (Socrates holds fast the objective uncertainty), and Høffding's own p. 23 makes the Socratic standpoint "objectiv Uvished". "Vished" is an error (printer's or copying slip) and is rendered by its evident sense, "uncertainty". The parenthesis is Høffding's gloss.
can be reached is an approximation, hence no qualitative decision, which comes only through faith. ``Instead of the objective uncertainty (on the Socratic standpoint), there is here the certainty that, viewed objectively, it is the absurd, and this absurd, held fast in the passion of inwardness\footnote{``Passion is precisely the tension in the contradiction; when this is taken away, the contradiction is a pleasantry, a bon mot''. Unsc. Postscr. p. 294.} is faith.''

Johannes Climacus has thus set up here a principle of faith that is incompatible with knowledge. And it is not only at the outset, when subjectivity chooses faith, that faith repels and is repelled by knowledge; what holds in principle for the development at its beginning must also hold in its further course.

This principle of his Johannes Climacus now sets against the two tendencies indicated earlier. Against Strauß and Feuerbach he maintains that science lacks competence to judge in the domain of existence; against the scientific doctrine of religion, that the bright shield of faith must not be stained with any spot of objective knowledge. That his polemic was directed chiefly against the latter was due in part, no doubt, to its being the prevailing one in this country, where negative criticism had no original representatives; but it was principally because, in his opinion, it mixed together modern culture and Christianity in an unclear and ambiguous way --- in this agreeing with the tendency of the time as a whole.

Here, then, we have the spectacle of three adversaries who make war on one another. If there could perhaps be talk of an alliance between two of them against the third, it would nonetheless obviously be only short-lived.

Attention naturally turns to the party that might perhaps seem to unite the other two in a ``higher unity'': dogmatics, which did, after all, maintain \emph{both} faith \emph{and} knowledge. It is clear that the presuppositions were now present for a
% --- p. 25 ---
\opage{25}\setcounter{footnote}{0}%
richer and deeper development of this science: on the one side, after all, the subjective principle of faith, on the other the objective principle of knowledge, had been carried through as sharply and consistently as can be conceived. Here was a victory to be won for whoever could master these opposites.

Then, in the year 1849, Martensen's ``Christian Dogmatics'' appeared: ``a scientific presentation and grounding of the Christian doctrines of faith in their inner coherence'' (§ 1).

% TR: "Enten-Eller" here is Mynster's disjunction, not Kierkegaard's book; rendered ``Either-Or'' in roman.
For anyone who remembered what demands the author himself had earlier made of a scientific dogmatics --- how, on the occasion of the ``Either-Or'' between supernaturalism and rationalism set up by Mynster, he had upheld mediation and declared\footnote{See ``Tidsskrift for Litteratur og Kritik'' 1839, especially p. 464, 469.} that a return to the Christian truth of revelation, except as a ``necessary turn lying within the philosophy of our time itself'', produced by carrying through Hegel's dialectical method, was unscientific (or at any rate ``no scientific advance'') --- for such a one there was every reason to examine the new ``Christian Dogmatics'' with particular regard to its scientific principle. Had that program been held fast, and had there even merely been an attempt at carrying it out, then things would have begun to look serious for negative criticism: it had, after all, been promised that in the ``immanent cognition of Christology'' the Straußian Christology too would be taken up as a ``relative mode of conception'', a ``sublated ideality''.

For anyone who had seen with how much steadfastness and persistence a whole series of pseudonymous authors, whose aim was to uphold existential truth and faith as its highest, paradoxical form, had hurled one challenge after another at ``the System'', which was, to be sure, not yet ``finished'', but was nonetheless expected soon to become so --- for such a one there was every reason to test how
% --- p. 26 ---
\opage{26}\setcounter{footnote}{0}%
the new ``Christian Dogmatics'' managed to reconcile the paradox with science, existential with scientific cognition.

``The question,'' it runs in the first of the authors who came forward against Martensen's dogmatics\footnote{``Mag. S. Kierkegaards `Johannes Climacus' og \textit{Dr.} H. Martensens `Christelige Dogmatik'. En undersøgende Anmeldelse af R. Nielsen''. 1849.}, ``around which theology and philosophy at the present time must be said to move, as around the true and proper question of principle, can be briefly stated in these words: Will the truth of Christianity by its nature be an object of objective knowledge?.. As a question of principle it can in no way be held off in evasive indeterminacy, but must be answered unconditionally, with Yes or No.'' The author then sets side by side the two different answers that the pseudonym Johannes Climacus and the author of the ``Christian Dogmatics'' give: the one with ``an unconditionally decisive No'', the other --- ``though, to be sure, with some uncertainty'' --- with Yes. The critique is then directed against this latter answer, in order to show that ``the `Christian Dogmatics' has undialectically circumvented the Christian problem'', whereas the pseudonym sharpens it. It is shown first that Martensen has not succeeded in separating \emph{dogmatics} and \emph{philosophy of religion}, ``theological'' and ``philosophical speculation'', but that on the contrary he contradicts himself (cf. § 2 of the Dogmatics with § 8, note). In order now to shed closer light on this unclarity, attention is turned to ``the problem of speculation''; it is asked whether the ``Christian Dogmatics'' acknowledges ``\emph{the demands that objective thinking imposes on the one who speculates}'', which are chiefly: ``a direct relation between theoretical certainty and speculative insight'', and, as regards method, ``the dialectical relation between the presupposition and the abstraction''. These demands would require that all dogmas be transformed into problems\footnote{%
% Run-over footnote: this note begins on p. 26 and continues on p. 27; the carry-over point is marked below.
Cf. ``Tidsskrift for Litteratur og Kritik'' 1839, p. 464, where the author of the Dogmatics declared that ``science could not fully
% [the note is carried over to p. 27 at this point]
be reconciled with any dogma before it was seen in its dialectical necessity in the System''.}. But that is something
% --- p. 27 ---
\opage{27}\setcounter{footnote}{0}%
the ``Christian Dogmatics'' cannot engage in: the absolute truth of Christianity is unshakable for it: it \emph{believes} in it. Therefore, in the third place, the critic turns his attention to ``the problem of faith'', and in particular to \emph{the relation between faith and its object}. According to the Dogmatics this relation is a double one: an existential and a theoretical relation. Very well --- says the critic --- then set up the objective \emph{object of faith}, baptism and the Lord's Supper e.g., and let us --- in faith --- obtain a scientific-objective decision as to what is the true doctrine of baptism or of the Supper! But see: the scientific theory depends on whether it is a \textit{theologia regenitorum} or \textit{irregenitorum} (rationalists), and, as regards the former, again on whether the regenerate person is a Lutheran (and here again: orthodox or pietist, ``Scripture theologian'' or Grundtvigian) or Reformed etc. So it is seen: how the objective object of faith is depends on faith, on the subjective. And if one then starts from \emph{faith}: how, then, can ``a cognition that is as subjective, as inward, as absolutely personal as a true cognition of faith must necessarily be, be converted into the objective, impersonal form of science and still be cognition of faith?'' On this the Dogmatics gave no enlightenment. Since, then, the ``Christian Dogmatics'' has \emph{neither} set up and upheld a distinctive theological principle of cognition, independent of philosophy, \emph{nor} acknowledged and carried through the principle of objective thinking, \emph{nor} resolutely and without half-measures placed itself under the banner of faith, or at least presented the principle of faith pure and simple --- the critique ends with the thesis: the ``Christian Dogmatics'' is without principle\footnote{What the critic's verdict was when he disregarded the principle, and it was merely ``a matter of judging this voluminous work solely with regard to the author's talent, learning and gift of presentation'': see the ``undersøgende Anmeldelse'' p. 42---44.}.

% --- p. 28 ---
\opage{28}\setcounter{footnote}{0}%
% TR: Danish "eu" (turned n) read as "en" (a). The guillemets »…« round the Latin tags on pp. 28--29 are kept as printed (TRANSLATION-PLAYBOOK §2: foreign terms in guillemets stay in guillemets); the Latin is kept as printed, including the spelling "qværo" (= quaero).
% TR: the Danish opens Stilling's title in this note and the inner Martensen title but closes only once; the English supplies the balanced pair.
The second author who came forward with a critique of Martensen's dogmatics sets out in his treatise\footnote{``Om den indbildte Forsoning af Tro og --- Viden, med særligt Hensyn til Prof. Martensens `christelige Dogmatik'.'' Af Mag. P. M. Stilling. 1850.} two different relations between knowledge and faith. The first is: »\emph{\textit{qværo intelligere ut credam}}«. Here, then, faith is not the presupposition. The spirit that strives after cognition wanders with its intellectual love through the various spheres of objectivity, and does not suspect, ``in the zeal of the process of cognition and in the first working days of the theoretically reflecting spirit'', that in the religious sphere it is situated otherwise than in the other spheres. ``In the religious sphere, too, the Hegelian determination of truth as process, synthesis of the subjective and the objective, holds good, to be sure . . . . But --- note well, and not to be forgotten --- \emph{in such a way} that the objective in this sphere is subjectivized, appropriated and made inward only in faith and hope, and by no means in a speculative-theoretical process''. It soon comes to light, therefore, that \textit{pro} and \textit{contra} here stand against each other --- that there are ``two ways and no certainty beforehand''. But before subjectivity resolves to seize the only certainty that can be in question here, that of faith and hope, serious theoretical and practical struggles will perhaps take place. ``And least of all could one who reflects, of `those valiant ranks', hit upon letting his power of thought be engaged as the yes-man of the dogmas; an engagement that is, to be sure, sometimes offered it by dogmatizing thinkers''. --- The proposition that the process of subjectivization in the religious sphere is qualitatively different from the theoretical one ``is forgotten in greater or lesser degree in Hegel, Marheineke, in short in all speculative theologians. It is of no use at all that Prof. Martensen says:
% TR: the Danish opens Martensen's words („En Troende...) but has only the one closing mark after "grundfalsk"; since "Det er grundfalsk" is Stilling's verdict, the English closes Martensen's words after "process." and supplies the balanced pair.
`A \emph{believer} can appropriate the Christian in a speculative-theoretical process.' That is fundamentally false''. To show this the author passes
% --- p. 29 ---
\opage{29}\setcounter{footnote}{0}%
on to the second standpoint: »\emph{\textit{credo ut intelligam}}«. This is ``reflection \textit{post fidem} (i.e., after I have taken the side of faith)''. But --- says the critic --- ``things go no better at all with reflection \textit{post fidem} than \textit{ante fidem} (i.e., before I took the side of faith and was still stuck right in the middle of the oppositions of reflection). The opposition of reflection (unbelief) was overcome in the first step of faith, not by reflection, but by and in believing living, which dictated obedience . . And every following step (2. 3. 4 . . .) is of the same qualitative nature as the \emph{first} step, which is continually repeated in the further course of the one who believes and hopes. And there may indeed be talk, and there is talk in Scripture too, of progress and growth in faith and hope, namely in the infinite strength of the subjective, personal certainty. But there is surely no talk of `growing beyond' in growth, or in `progress' going beyond the qualitative limits of the certainty of faith and hope to the \emph{objective} certainty of proof''\footnote{Op. cit. p. 37 f.}. Since, however, the ``Christian Dogmatics'' had set up precisely such a »\textit{credo ut intelligam}« as its principle, the critic asks: ``Has Prof. Martensen really succeeded in lifting the cross of thought and in obtaining from the power of reflection its sincere support and sanction of the dogma?'' And the result at which the critic then arrives, after an examination of the Dogmatics' treatment of the most important dogmas, is ``that »\textit{contra}« has not been vanquished at all, but only ignored or hushed up; that the testimony of reflection has been heard only with partiality, and is never allowed to have its full say'', ``that our dogmatician's »\textit{credo}« is in practice only a scholastically inhibiting \textit{credo}, which places the freedom of reflection under censorship and permits it no utterances other than such as are well-pleasing to the dogmatician's ear'' (p. 66). ---
% TR: Danish »contra“ (guillemet opening, high comma closing) rendered as a balanced »contra«.
And the critic finds something teleological in the fact that »\textit{contra}« cannot be scientifically set aside. For if there were a ``proof \textit{post fidem}'', this would have an unfortunate retroactive effect
% TR: p. 29 ends mid-sentence after "Tilbagevirkning", rendered "an unfortunate retroactive effect"; its object ("paa Troes- og Haabsvisheden", p. 30) is left to the next batch.
% ==== Batch 4: printed pp. 30--38 (PDF 34--42). English translation of pot/src_b4.tex.
% --- p. 30 ---
\opage{30}\setcounter{footnote}{0}%
% TR: the slice opens mid-sentence: p. 29 ends "Gaves der nemlig et „Beviis post fidem“, saa vilde dette have en uheldig Tilbagevirkning" and p. 30 goes on "paa Troes- og Haabsvisheden."; rendered "...then this would have an unfortunate retroactive effect" | "on the certainty of faith and hope." No \ldots, since the preceding page is translated.
on the certainty of faith and hope. ``For it is psychologically
impossible, indeed only a pompous phrase, that I should at once
have the certainty of faith and hope ``independently of speculation,''
``by an entirely different way than speculation,'' and yet in the
same breath, precisely by the ways of speculation, get the certainty
of faith and hope more and more (the boundary grows fluid . . .)
confirmed'' (p. 64 note).

It cannot be denied that the critique had stated clearly and
definitely what the point at issue was. What the objections
were directed against was not this or that particular, but
the dogmatic principle itself. The demand therefore necessarily
had to be made of the anti-critique, if one should appear, that
it should not --- dogmatically --- regard the dogmatic principle as
standing unshakably firm, but should enter upon a closer
grounding of it, and should develop clearly and unequivocally its
relation to ``pure'' knowledge on the one side, and to faith
on the other.

% TR: „Dogmatiske Oplysninger“ (Martensen, 1850; "Dogmatic Elucidations") is a title and stays Danish throughout; the author's play on it (the "Oplysninger" that must "oplyse", elucidate) is lost. Dogmatiken = "the Dogmatics" where Martensen's book is meant, "dogmatics" for the discipline.
But the answer\footnote{``Dogmatiske Oplysninger''. An occasional piece by \textit{Dr.} H. Martensen. 1850.} that the author of the Dogmatics gave his
critics was now precisely a purely dogmatic one. It is indeed ``an
old experience that dogmatics always has critique in its
train,'' and ``especially in our time, which is far more critical than
dogmatic''; but as regards the critique of the present Dogmatics,
the author thinks that it ``cannot yet be said to
have begun'' (p. 7). This might seem strange to one
who has read the two critical treatises mentioned above; but
on closer examination one discovers that the critique is rejected
because it attacks the principle itself: when it does that, then
from the dogmatic standpoint it is not ``proper'' critique, indeed
is really not noticed at all by dogmatics. Thus it
takes up its position, e.g., toward Strauß, and toward
% TR: Danish „euhver“ (turned n) rendered as „enhver“, "any".
any ``thoroughgoing doubt about the foundation of Christianity''; however
% --- p. 31 ---
\opage{31}\setcounter{footnote}{0}%
``significant'' such a doubt may otherwise be in other respects:
it does not concern dogmatics at all! And why not?
Because ``the dogmatic interest'' does not exist at all for such
a doubter!\footnote{See Martensen's Dogmatics § 31 note.}

Here, then, a secure barrier has been set against all critique.
There could thus be no question of any real exchange
between dogmatics and critique. Otherwise, however, the
% TR: „dogmatiske Oplysning“ (singular, lower case) here plays on the title; rendered ``dogmatic elucidation''.
``dogmatic elucidation'' varies in character.

Toward Mag. Stilling, e.g., the ``dogmatic'' showed itself
% TR: „Hr. Magisterens Skriveri“: Hr. rendered Mr. by the brief's rule, hence ``Mr. Magister's scribbling''.
in this, that ``Mr. Magister's scribbling'' was at once ``laid aside, not
to be taken up again'' (p. 8). Mag. Stilling rightly remarked
on this in a later treatise: ``Of the invincible and
most honorable weapons with which a speculative dogmatician is able to
maintain his standpoint against thoroughly argued attacks
from the side of former friends, Prof. Martensen's
anathema upon my critique affords a very instructive example.''

The weightiest critique of the Dogmatics had already been written
before it appeared: by the pseudonyms. Nevertheless, in
the Dogmatics itself no regard at all was paid to these authors,
except that in the preface ``some among
us'' are reproached for wanting to be ``absolute in faith''\footnote{As one easily sees, an ambiguous expression: for it is one thing to maintain
the principle of faith absolutely, without abating it, and another to declare
oneself to stand at the absolute height of faith. --- For the rest,
the preface also alluded to ``a few individuals who, without feeling
any urge toward coherent thinking, can satisfy themselves by
thinking in stray thoughts and aphorisms, notions and flashes.'' Thus
ran the verdict in 1849. Would it also run thus in 1866?}. Now Prof.
Nielsen in his critique had brought Johannes Climacus forward again, so
that what had been neglected might perhaps be made good. The ``Dogmatiske
Oplysninger'' consequently had to elucidate the relation of the Dogmatics to
% TR: Danish „Pseydonymerne“ (y for u) rendered as „Pseudonymerne“, "the pseudonyms".
the pseudonyms. Now although these authors are declared ``wholly
irrelevant in the present connection,'' although the author's
% --- p. 32 ---
\opage{32}\setcounter{footnote}{0}%
``acquaintance with this extensive literature is only very slight
and fragmentary,'' he nevertheless declares that he has an entirely
different conception of these writings than the reviewer: namely, that
they did not mean to ``ground a reform of dogmatics'' but ``in
Socratic fashion to set a deeper skepticism in motion, in order
thereby to awaken the reader to seek the problem that is
higher and more than all theological and philosophical school problems,
namely the personal problem of life'' (p. 12). And
one need not have read much in the pseudonymous writings
to confirm the truth of these words. But, in the first place, the author of
the Dogmatics was not at all at odds with his reviewer here; for the latter
had indeed precisely wanted to prevent his ``conception of the book's
% TR: „Uvidenskabelig Efterskrift“ (and „Uvidensk. Eftersk.“ in note 2 below): the author's shortening of the Concluding Unscientific Postscript, rendered ``Unscientific Postscript'' / "Unscient. Postscr.".
(i.e., the ``Unscientific Postscript's'') scientific significance from being confused
with the author's perhaps somewhat divergent\footnote{With this may be compared the relation, so characteristic of Kierkegaard,
in which he placed himself to this polemic. See ``Kierkegaards
Bladartikler'' p. 252 note.}
% TR: „undersøgende Anmeldelse“ (words from the title of Nielsen's 1849 review, "examining review") kept in Danish as a title; the print has no closing quotation mark after it, supplied here.
main thought'' (Prof. Nielsen's ``undersøgende Anmeldelse'' p. 7).
And secondly --- must not the ``science of faith'' fall if
Johannes Climacus's concept of faith is correct? Unless,
that is, one regards the ``deeper skepticism'' he wants to set in
motion as so deep that it acquires no real significance.
--- For the rest, the ``Dogmatiske Oplysninger'' will
have nothing to do with Joh. Climacus himself; he is dismissed as
irrelevant; but only with Prof. Nielsen's presentation of
his --- ``theses.'' It is then declared that ``what is new in them
is not true, what is true not new''; thereupon a number of
old truths are discussed --- without any inquiry into whether the new
might not lie precisely in the manner in which they were set forth\footnote{``Objectively the accent falls on: \emph{what} is said; subjectively: \emph{how} it
is said.'' Unscient. Postscr. p. 149.};
until the author stops at two ``new theses'' in order to prove
that they are not truths. The first of these, ``the truth,
inwardness and subjectivity,'' was new to the author ``in the
% --- p. 33 ---
\opage{33}\setcounter{footnote}{0}%
sense that he had not found it in any dogmatics known to him''
(p. 16). This was not so strange after all, since
the thesis (at least in its secondary thought, its consequence) is
polemical against scientific dogmatics. The other thesis
that was new to the author was: ``for the existing subjectivity
the highest truth is the paradox.'' In this thesis the
author rightly thought that ``the problem was to be found''; but
unfortunately he forgot that in order to recognize that the truth,
the highest truth, is the paradox, one must necessarily proceed
from the standpoint of the \emph{existing} subjectivity: what
wonder, then, that his objective-dogmatic investigation arrives at
the result that ``paradoxical'' is only a rhetorical hyperbole, and
the paradox something invented by the serpent in Paradise
(p. 23)!

% TR: Erkjendelse = cognition (repo term) here and below.
If one asks what the ``Dogmatiske Oplysninger'' teach
about the relation between existential and scientific cognition,
there are two points in particular to which we may draw
attention. In order to show that faith contains within itself
a germ of science, the author insists that faith is not merely
``consciousness of redemption'' but ``consciousness of revelation'' (p.
30--40). Failure to heed this is supposed to be the basic error
in Johannes Climacus's concept of faith, as set forth by Prof. Nielsen.
But what in the world does this elucidation help? Does
faith become more scientific because it is consciousness of revelation?
Or is not revelation precisely true only for the \emph{believer},
and only \emph{insofar as} he is a believer? Is it a \emph{scientific}
striving ``to seek revelation in redemption,''
``by the way of devotion and believing contemplation to seek to
know the God who through his work of redemption discloses
his essence to our gaze''? But let us nevertheless assume for a moment
that dogmatics succeeds in getting started on its
task: ``scientifically to develop the contemplative moment of faith,''
--- of what kind, then, will the scientific
cognition be that results from it? Here the author maintains --- and
% --- p. 34 ---
\opage{34}\setcounter{footnote}{0}%
this is the second point we wish to bring forward --- that ``dogmatic
inquiry into and searching of the truths of revelation''
is not ``a variety or a corruption of philosophical
speculation,'' but has ``its own principle, independent of philosophy,
and should therefore also be judged by its own standard''
(p. 66--68); for ``the redemption of fallen human nature''
must also embrace the redemption of thought, as surely as it
is ``the purpose of Christianity to restore true
humanity completely'' (p. 22).

There is thus a specifically theological-scientific principle of cognition
--- which is to be ``nothing other than what
lies in the rightly understood doctrine of the \textit{testimonium spiritus
sancti}'' (Dogmatics § 29). This is precisely a crux
of the whole dispute. Prof. Nielsen had already directed his attention
to it in the ``undersøgende Anmeldelse''; in the
% TR: „Belysning“ (Nielsen's reply, whose title reads „... Dogmatiske Oplysninger belyste“) rendered as a description, ``Illumination''; the book title in the note stays Danish.
``Illumination'' of the ``Dogmatiske Oplysninger'' with which he
concluded the polemic\footnote{``\textit{Dr.} H. Martensens ``Dogmatiske Oplysninger'' belyste af Prof. R.
Nielsen''. 1850.}, it is said with regard to the aforementioned
``theoretical significance of the \textit{testimonium spiritus sancti}'': ``According
to \textit{Dr.} Martensen's ``basic view,'' then, there are to be
two kinds of objective cognition, namely: a lower, worldly,
general-scientific one, which cannot yet say: \textit{Credo ut intelligam},
since it is still in its sinfulness, . . . and a
higher, ecclesiastical, dogmatic-scientific cognition, which can straightway
say: \textit{Credo ut intelligam}, since it is redeemed by the grace
of God, and so redeemed that it can produce the genuine
theological speculation . . . . In the speculative pulpit style,
which preaches speculation without thinking the thought, such
an assurance about ``the redemption of thought'' may certainly look
quite good, both upbuilding and speculative; but --- ``show
me thy faith by thy\footnote{Cf. Grundideernes Logik I, p. 44--53, where it is shown in a single example
how little the deed here answers to the faith.}
works!''\,'' (p. 37--38).

% --- p. 35 ---
\opage{35}\setcounter{footnote}{0}%
For one must ask: does one have a different logic and
ontology on the dogmatic standpoint (i.e., after having
been regenerated) than otherwise in science? Now since it is
clear that theology in this respect is subject to the same conditions
as the universally human sciences, where then is
this logic and ontology to be checked? Surely in philosophy,
which ``by its logical and ontological investigations
forms a propaedeutic foundation for all science'' (Martensen's
Dogmatics § 36)? But how then can the ``non-believing,'' ``non-regenerate''
science check the ``believing'' and ``regenerate''?
And what other outcome will this checking have than that
the ``\emph{believing} science'' must be assigned ``a place of its own
\emph{outside} science''? And if one then wished to answer this
by citing 1 Cor. 2, 14--15, that would at once
be a confirmation that this verdict was right, and a
great misunderstanding of what ``the Spirit of God'' means.

A principle, to be sure, is not to seek its grounding outside
itself, but to be something self-valid. But a scientific principle
must be able to give a scientific account of itself. It is of no use,
when the checking comes too close, then to appeal to
faith and declare that on the basis of it one has a higher
scientific standing.

And herewith, then, we conclude our account of this polemic.
It has often been said that such literary polemics are
unfruitful and without result. And certainly it may be difficult
to state in a finite sense the result at which one
has arrived. But by every able polemic this at least is achieved,
that the contending principles get a hearing and step
forth onto the field of battle to venture a joust with one another. And
will it not then soon become manifest who,
ideally regarded, is the vanquished? When objections have been raised
that strike the matter in its innermost essence, and one
then simply dismisses them without properly engaging with them, is
that not a defeat?

% --- p. 36 ---
\opage{36}\setcounter{footnote}{0}%
And precisely in this manner did dogmatics, in the
polemic treated here, conduct itself toward critique. That I think I
have shown by the foregoing account. If one wishes to call it
partisan and one-sided, I answer: show me another way
of understanding the matter; show that dogmatics has really engaged
with critique and overcome it! show that dogmatics has
been able to appear with a good conscience before the tribunal of science
as well as before that of faith, that it has been able to unite both
without halving them and abating both!

The three combatants who earlier stood on the field of battle
against one another have thus been reduced to two. But
this duel gives occasion to several misgivings.
On the one side stands faith, on the other side --- well,
what really stands there now? Strauß asserts the modern consciousness
and its world-view in opposition to the standpoint of ``the old,
especially oriental world,'' from which the biblical
writings were composed. ``The modern age has a series of
laborious investigations, continued through centuries,
to thank for the insight that everything in the world hangs
together by a chain of causes and effects that tolerates no
interruption . . . . From this standpoint the
biblical narratives cannot possibly appear as history . . . . The
divine cannot have happened thus; what happened thus
cannot have been divine.'' How, then, does the modern,
scientific consciousness here take its position toward Christianity?
Well, it takes seriously that the researcher of religion is to
relate himself to revelation as the natural scientist to nature:
namely, in such a way that he seeks in a rational, scientific manner
the laws that lie at the ground. This standpoint
Hegel had approved with respect to mythology\footnote{``Die Hauptsache der Mythologie ist das Werk der phantasirenden
Vernunft, die sich das Wesen zum Gegenstande macht, aber noch
kein andres Organ hat als die sinnliche Vorstellungsweise; daher die
Götter als Menschen erscheinen. Die Mythologie
% TR: German „kan“ (for „kann“) kept as printed.
kan nun studirt
werden für die Kunst u. s. w.; der denkende Geist muß aber den substantiellen
Inhalt, den Gedanken, das Philosophem, das implicite
% TR: German „inthalten“ (for „enthalten“) kept as printed.
darin inthalten ist, aufsuchen, \emph{wie man in der Natur Vernunft
sucht}''. Geschichte der Philosophie. I, p. 98.}; it was Feuerbach who
% --- p. 37 ---
\opage{37}\setcounter{footnote}{0}%
with great acuteness carried this further and sought to find
the psychological laws that had guided the origin and development
of religion, and especially of Christianity\footnote{``Das Wesen des Christenthums''. Von Ludwig Feuerbach. Leipzig.
1841.}. Thus he came
to the thesis: ``Religion ist das bewußtlose Selbstbewußtsein
des Menschen'' (p. 37). Before human self-consciousness
is properly developed, it cannot intuit its own universal
essence, the idea as idea, but clothes it in a form of fantasy,
personifies it, and endows the personal God thus arisen
with human attributes. Religion has
thus arisen through one's having made subject what
should be predicate, and vice versa; thereby the illusion of
an opposition between the divine and the human has
arisen. ``Gott ist der Begriff der Gattung als Individuum''
(p. 202). ``Das Bewußtsein des Menschen in seiner
Totalität ist das Bewußtsein der Trinität'' (p. 72). Theology,
therefore, is anthropology. What must now be done,
now that reason and self-consciousness have come of age, is:
to make what for religion is the main thing, the subject,
into the subordinate thing, the predicate. ``Das Richtige, Wahre und Gute
hat überall seinen Heiligungsgrund in sich selbst, in seiner
Qualität. Wo es Ernst mit der Ethik ist, da gilt sie eben an
und für sich selbst für eine göttliche Macht'' (p. 375 f.).

Strauß and Feuerbach say that they speak in the name of science.
But can they now also really count as the rightful representatives
of science over against religion? or must it
not be said to be a one-sided conception of scientific
reflection that it should testify only \textit{contra}, just as it
% --- p. 38 ---
\opage{38}\setcounter{footnote}{0}%
is one-sided when theology can hear only the testimony \textit{pro}?
And finally, would it not be a precarious relation in which
science came to stand to humanity and the spiritual
life, if over against religion it had to be absolutely destructive?

How the matter stands from the side of faith: that faith
has its own sphere, where it does not acknowledge the supremacy of science
(and were it once to engage with objective reflection,
it would be in a bad way: for here there are, after all, both \textit{pro} and
\textit{contra}, two ways!), --- this we have already seen in the
foregoing. It is S. Kierkegaard who has fought for this,
the religious side of the matter.

But science, philosophy, too, must of course become clear with
itself about what its right relation to religion is. From
this side Prof. R. Nielsen has worked on the problems
belonging here. And it is clear that the settling of this point
must be brought about by the scientific route.

The first question, then, is: why do the results of science
always sound \textit{contra} revelation for Strauß and Feuerbach?
Presumably for the same reason by which theology always hears
a \textit{pro} from the mouth of science: the freethinker's science is an
unbelieving science, theology a believing science. But faith
and unbelief are not scientific principles: they belong to life,
to existence.

Science asks only about the definite, objective grounds.
``From the standpoint of life, on the other hand, subjectivity reigns''; here,
then, everything is decided by subjective grounds. Now when, therefore, a
decision of questions that belong to the sphere of existence, of the personal
and religious life, gives itself out to be
scientific, it is clear that the two domains are confused,
--- whether the decision then falls out in the one direction or the other.
``The Straußian endeavor to raise science to
religion cannot possibly be heterogeneous with the theological
endeavor to raise religion to science; for both
% TR: the slice ends inside a word: p. 38 "thi begge Be-" | p. 39 "stræbelser gaae jo ud fra den Forudsætning, at Religion og Videnskab ere eensartede“ (Evangelietroen og Theologien p. 82)". The English breaks at "for both" and leaves the whole word "endeavors" ("for both | endeavors proceed, after all, from the presupposition that religion and science are homogeneous''") to the p. 39 side; no \ldots.
% ==== Batch 5: printed pp. 39--46, end of the book ====
% Page 39 begins mid-word: p. 38 ends "Be-", the word being Be|stræbelser.
% --- p. 39 ---
\opage{39}\setcounter{footnote}{0}%
% TR: the sentence runs over from p. 38: "Den straußiske Bestræbelse, at
% TR: hæve Videnskaben til Religion, kan umulig være ueensartet med den
% TR: theologiske Bestræbelse at hæve Religionen til Videnskab; thi begge
% TR: Bestræbelser gaae jo ud fra den Forudsætning, at Religion og Videnskab
% TR: ere eensartede." This slice renders "[Be]stræbelser gaae jo ud fra ...
% TR: eensartede" as "endeavors proceed, after all, from the presupposition
% TR: that religion and science are homogeneous"; the opening quotation mark
% TR: stands on p. 38.
endeavors proceed, after all, from the presupposition that religion and
science are homogeneous'' (Evangelietroen og Theologien p. 82).

There is, then, nothing else for science to do but to
pronounce its \textit{non liqvet} concerning these questions and to leave
their decision to the sphere where they belong. It
becomes merely comical when it comes forward here with its objective
reasons and fills the scales, the one with \textit{pros}, the
% TR: "pro'er", "contra'er": the Latin words with a Danish plural ending,
% TR: in antiqua; rendered "pros", "contras", kept in \textit.
other with \textit{contras} --- and behold! after all it does not depend in the least
on these reasons to which side the scale sinks.

``But,'' Strauß will say, ``everything in the world is, after all, connected
by a chain of causes and effects; the miracle must therefore be
rejected by science as a sheer impossibility.'' To this the
answer from the side of philosophy is that this must certainly
seem so to the physicist; for ``the causes of physics are causes of things,
its laws are laws of necessity; but the highest cause of philosophy
is an idea-cause, and the idea-cause a teleological
self-cause. Inasmuch as the absolute self-cause lays the laws of necessity
at the foundation of the law of freedom, there is seen from the height of
idea-cognition a depth of possibilities which no human
knowledge is able to measure out. Instead of fantasizing
about the indeterminable and brooding over the unfathomable,
philosophy here holds fast to the essential task:
to secure the right of science by stating the limit of knowledge''%
\footnote{Grundideernes Logik. I, p. XXIV.}.
Philosophy, then, could not possibly assert the validity of spirit and of the
ideal in existence if it were to \emph{break} with
religion, where the absolute omnipotence of spirit is asserted. If this holds
in the intellectual domain, then it holds even more
in the ethical. For it is not the case, as Feuerbach
thinks, that ``das Richtige, Wahre und Gute'' would be able to hold
its own in actual life if it were not supported by religion.
The idea of the good would soon be eliminated by human
% --- p. 40 ---
\opage{40}\setcounter{footnote}{0}%
reflection%
\footnote{This is developed in detail in ``Forelæsninger over philos. Propædeutik
1860--61''. p. 416--423.}. As surely, therefore, as philosophy is to
have an ethics, where ``the thoughts of life meet with the scientific
fundamental thoughts,'' so surely is it true that it dare not declare the Gospel
to be a myth (Grundideernes Logik XXVII).

But philosophy can just as little comprehend and ground the Gospel.
``In declaring that objectivity is
the truth, science by that very act renounces all influence
on faith.''

This does not, however, exclude that religion and
personal life in general can become an object of science.
But this can then only happen in such a way that the subjective
is ``objectified in reflection,'' and the religious view is thus
had ``at second hand,'' ``not in a religious way,'' but
``through the medium of thought and imagination'' (cf. Martensen's
Dogmatics § 8, note). Here one must then hold fast to the
categorical ``difference between the scientific in the wider and in the narrower sense''
(R. Nielsen's philos. Propædeutik p. 78).
Scientific only in the wider sense is every treatment of
an object that lies outside the sphere where science
rules with absolute \textit{dominium}; hence every presentation
of the subjective, personal life. What is scientific
lies in the objective treatment. Such a ``science of subjectivity''
can then have the religious truths as its object,
not in order to ground them, but in order to ``illuminate them as many-sidedly
as possible,'' in that it ``conceives the ideals of faith as ideals of life
and finds in the end aimed at the standard for judging
the means'' (Grundideernes Logik XXVI).

Now if it is the case that there is a qualitative
difference between faith and knowledge, that they belong to entirely different
spheres, then the possibility of their
reconciliation is thereby given. Where, in what domain, in what sphere
% --- p. 41 ---
\opage{41}\setcounter{footnote}{0}%
is this reconciliation to take place? In the sphere of knowledge? But
then knowledge itself is both party and judge; the mediation will
therefore be partial, the unity one-sided (cf. Unscient. Postscr.
p. 286 f.). In the sphere of faith? Here the answer can be Yes,
if by ``the sphere of faith'' one understands the sphere of personal
existing as a whole. It is the existing subjectivity
that has a relation both to knowledge and to faith, in
whose consciousness the conflict is waged. ``The conflict is psychological; it
is human frailty, which has neither resignation
enough to reconcile itself with what is objective in science, nor inwardness
enough to hold fast to faith'' (Philos. Propæd. p. 77).

If the heterogeneity between faith and knowledge were not
absolute but only relative, so that both belonged to the same sphere,
how would they then be able to be united in one consciousness without
contradiction? When faith says one thing and science something entirely
other and opposite, how would one then be able to reconcile these
conflicting assertions? If, on the other hand, it is recognized that one thing is existential
and religious truth, another scientific truth, then the rule can
% TR: the print has »Suum cuiqve!» (antiqua guillemets, both printed as »);
% TR: the English gives the balanced pair »...«. "cuiqve" kept as printed.
hold: »\textit{Suum cuiqve!}« Then it is clear that the psychological
reconciliation \emph{can} (for whether it \emph{actually} comes about
depends, after all, on the single existing subjectivity) take place,
namely in such a way that faith is set highest, while science
is nevertheless retained as a relative task; for the reverse is
not possible%
\footnote{Thus it is precisely \emph{because} faith and knowledge are absolutely heterogeneous that their
psychological reconciliation is possible. It is a great misunderstanding to put, instead of
this ``because,'' a ``\emph{though}'' (see Mr. \textit{cand. mag.} Brandes's
article in ``Fædrelandet'' of January 13th). Then one can easily
find ``obscurities''!}.

The reconciliation between faith and knowledge, then, is not settled
dogmatically in the Nielsenian philosophy. On the contrary, everyone receives
a direction and a summons to accomplish the reconciliation practically and ethically
in the course of his own personal development of life.
What philosophy can teach us is to hold the principles fast,
% --- p. 42 ---
\opage{42}\setcounter{footnote}{0}%
not to distort and confound them. If it should then seem as
though the conflict could not be settled, as though the struggle never came to an end ---
then Ørsted's words hold: ``Let us give truth the honor!
With it the good is indissolubly bound up. The full truth
brings its own comfort with it.'' And the comfort is this, that one
has not denied the ideal.

It must also be remembered here that science is only a
single side, albeit an exceedingly significant one, of human
life. If the question, therefore, is the relation between Christianity
and science, this is only a particular form of
% TR: Humanitet / det Humane = humanity / the humane (the purely human
% TR: sphere of culture, art, family and civic life), at first occurrence.
the question of the relation between Christianity and humanity
as a whole. Science stands beside art, family life
and civic life, in short: everything that belongs to the
purely human, the relative over against what in Christianity
is set up as the highest end: ``the Kingdom of God.''

In these other directions, too, a confounding
of the Christian and the humane is possible; indeed it has taken
place, and it surely takes place still. Here, too, the point is
not to be finished with the ``reconciliation'' too early, but
to fix the oppositions in view as sharply as possible --- and
then to take one's side.

In Holy Scripture one finds no direct information
about the relation between faith and knowledge, between Christianity
and humanity. This has its ground in the fact that Christ
and the apostles stood at such a height of existence, so
absolutely at the religious standpoint, that the --- in relation
to this, relative --- realm of humane ideas became something altogether vanishing.
But this holds only for one who stands at the standpoint of
inspiration, where the Holy Spirit works unconditionally
and immediately, without engaging with and entering into interaction
with the given conditions (cf. Acts of the Apostles
2, 3--4).

In the time that followed, human culture developed
in interaction with, indeed for a long time under the guardianship of,
% --- p. 43 ---
\opage{43}\setcounter{footnote}{0}%
religion. But the consciousness of an opposition between
the two keeps coming back. If the opposition now becomes a
dualism, then only two ways are possible: the ascetic's or
the freethinker's. The one breaks with humanity, the other
with Christianity.

But let us look more closely: is the ascetic's standpoint
that of the faith of the Gospel, of original Christianity? One need
not have read much in the New Testament to know that
a violent and painful break between Christianity and humanity,
such as the one we find in the ascetic's life, is not found in
the apostle. He unites, in the power and loftiness of the spirit, the decisive
oppositions, without one noticing any pain,
any restless unclarity in him; whereas the ascetic writhes
as under a heavy burden.

Thus: ``Respect for the monastic movement of the Middle Ages!''
(Joh. Climacus) Respect for one who can bear the unconditioned
unconditionally in this world of conditions! But the ascetic
exemplar dare not obscure the evangelical one. The Gospel has
laid an interdict neither on science nor on the state nor on any
universally human condition.

From this it is concluded that it must be possible to unite the
Christian and the humane, that a dualism cannot possibly be
the right relation.

If Christianity really still exists, if it has not been
overcome by advancing culture, then it must be possible
to hold fast to the evangelical ideal in the midst of the humane world.
But it is just as certain that one must then acknowledge that one
lives under other spiritual conditions than the apostolic and New Testament ones.

For above all: no confounding. And a
confounding takes place where one does not heed the
essential and decisive opposition. Then one gets neither
true Christianity nor true humanity. ---

If the reconciliation of faith and knowledge is to be psychological,
existential, then it can be seen how little to the point the objection is that
% --- p. 44 ---
\opage{44}\setcounter{footnote}{0}%
by referring faith and knowledge to qualitatively different
spheres one secures one's standpoint of faith against all freethinkers'
objections: ``what risk is there in believing `against the concept'
when we have secured for ourselves in advance the insight that the concept has
no validity whatever in the sphere of faith?'' (``Dogmatiske
Oplysninger'' p. 74). For because it is recognized \emph{scientifically}
that knowledge has no authority in the domain of faith, does
it thereby become easier to \emph{believe}? On the contrary, the insight that
has here been won in \emph{knowledge} must be won again from the very beginning and on other terms
when it is a matter of \emph{believing}. For ``it cannot be
set up as a new dogma about dogma that dogma is a
% TR: Stilling's "Tankekors" (a cross/crux for thought) rendered "a crux for thought".
crux for thought. Only in the attempt to lift dogma does
the power of reflection make this discovery''%
% TR: the dash in the title ("af Tro og --- Viden") stands in the print; kept.
\footnote{Stilling: ``Om den indbildte Forsoning af Tro og --- Viden''.
p. 55 note.}. And as regards
the talk of ``the simple believers who without such a
preparation take their stand at the standpoint of faith,'' it must be remarked,
first: that it is, as has been said, a misunderstanding when
there is talk of ``such a preparation,'' and next: that here
the same ``simple'' difference holds between ``the wise man''
and ``the simple man'' as in all ethical cognition: ``the simple man
knows it; the wise man knows that he knows it'' (Joh.
Climacus). For the rest, it is hardly from the side of speculative
theology that there should be too much talk of ``the simple
believers.''

While here one has sought to present the view of
a difference of spheres between faith and knowledge as a new ``scholasticism,''
there are others who accuse it of scientific
``skepticism.'' It ``strikes'' Mr. Larsen (in ``Dagbladet''
of January 19th)%
% TR: "lyse i Ban" (to declare under the ban, excommunicate) rendered
% TR: "to pronounce the ban".
\footnote{From this article one sees, moreover, that if Mr. Larsen has learned nothing else from
the earlier conflict over this matter, he has at least learned to
pronounce the ban.}, he ``can in no way get it into
his head,'' how one can get a science that recognizes
% --- p. 45 ---
\opage{45}\setcounter{footnote}{0}%
``eternal laws of nature'' without directly declaring the miracle
impossible. He then states ``three possible standpoints'' toward
the miracle, but to little use for the reader; for one is not told
whether all three of these standpoints are scientific, or
whether one or two of them are, or whether perhaps none at all of
them is. And yet it was on this that the whole matter depended. On this
he should have consulted his ``expert building commission.''
From the same quarter he might perhaps also find out for me what
``strikes'' me, what I ``cannot get into my head,''
namely: how does one go about drawing the miracle
into science and alongside this holding on to natural science
with its ``eternal laws of nature''?

The same question that is treated so clumsily here
comes forward in a far clearer way in Mr. G.
Brandes's articles in ``Fædrelandet'' of January 13th and 20th.
I understand the question thus: how can philosophy,
in its own domain, scientifically justify the
regard it pays to a sphere foreign to it? --- With this
the discussion is led away from the boundary dispute between philosophy
and theology to the purely philosophical domain%
\footnote{Mr. Brandes stands in a peculiar relation to Mr. Larsen. At first
glance one would think that he was defending Mr. Larsen; but then he declares
the latter's defense of theology to be insignificant and
ascribes merit to him only as an attacker; and behold, then the point
where Mr. Brandes thinks the attack must be made is an entirely different one
from the one Mr. L. has attacked. Consequently Mr. L. retains only
the merit of having attacked Prof. Nielsen at all.}.
To that extent it does not properly concern the present treatise.
Nevertheless I will remark the following.

What is demanded can only be the explicit development
of what it means that ``the highest cause of philosophy'' is
``the absolute self-cause'' (Grundideernes Logik XXIV).
This development can, in Mr. Brandes's opinion, take place only
in the transition from the circle of knowledge to that of power and, in a still
% --- p. 46 ---
\opage{46}\setcounter{footnote}{0}%
more concrete form, in the transition from logic to the philosophy of nature.
But if Mr. Brandes really has found no
elements at all for the solution of his problem in what has so far appeared of
``Grundideernes Logik,'' then he should not have confined
himself to wishing for elucidation and completion, but
% TR: "Subjectivitetens Logik", "Functionernes Logik" (the logic of
% TR: subjectivity, of the functions): parts of Nielsen's work, kept in Danish.
then he would have had to find fault already in the development of ``Subjectivitetens
Logik'' and ``Functionernes Logik.'' For that lack
necessarily presupposes this fault, if indeed what has appeared of
``Grundideernes Logik'' really is a whole in itself and not a
fragment.

Besides, it surprises me that Mr. Brandes has not taken
account of the significance that is ascribed to ethics in the introduction to ``Grundideernes
Logik'' (p. XXVII). For if there is to
be a philosophical ethics, then it cannot stand in contradiction
with logic and the philosophy of nature. And since ethics
forms a transition between philosophy and religion, the
way is thereby shown by which one, without ``making science
religious'' or ``writing a new theology,'' can respect the
religious by recognizing the limit of science. ---

Here there may be much still to do. Nor has it been
promised that all questions shall be able to find their absolute
solution. ``It is not the case with the progress of our knowledge
that the remaining problems should become fewer, easier
and more insignificant the nearer we approach the ideal; on the contrary! the
more clearly the ideal is seen, the higher the problems become and the
deeper runs the stream of inquiries''%
\footnote{Grundideernes Logik p. 253.}. --- But what
can be attained is that essential progress is made, that one
comes to be placed more rightly. And herein lies, in my conviction,
the great significance of the view in the philosophy of religion
presented here, whose historical presuppositions, development and
essential content I have endeavored to make clear
in this little treatise.

% Closing rule, centred, below the footnote on p. 46, as in the print.
\begin{center}\rule{0.2\textwidth}{0.8pt}\end{center}

\end{document}
